% Options for packages loaded elsewhere
\PassOptionsToPackage{unicode}{hyperref}
\PassOptionsToPackage{hyphens}{url}
\documentclass[
]{article}
\usepackage{xcolor}
\usepackage[margin=1in]{geometry}
\usepackage{amsmath,amssymb}
\setcounter{secnumdepth}{-\maxdimen} % remove section numbering
\usepackage{iftex}
\ifPDFTeX
  \usepackage[T1]{fontenc}
  \usepackage[utf8]{inputenc}
  \usepackage{textcomp} % provide euro and other symbols
\else % if luatex or xetex
  \usepackage{unicode-math} % this also loads fontspec
  \defaultfontfeatures{Scale=MatchLowercase}
  \defaultfontfeatures[\rmfamily]{Ligatures=TeX,Scale=1}
\fi
\usepackage{lmodern}
\ifPDFTeX\else
  % xetex/luatex font selection
\fi
% Use upquote if available, for straight quotes in verbatim environments
\IfFileExists{upquote.sty}{\usepackage{upquote}}{}
\IfFileExists{microtype.sty}{% use microtype if available
  \usepackage[]{microtype}
  \UseMicrotypeSet[protrusion]{basicmath} % disable protrusion for tt fonts
}{}
\makeatletter
\@ifundefined{KOMAClassName}{% if non-KOMA class
  \IfFileExists{parskip.sty}{%
    \usepackage{parskip}
  }{% else
    \setlength{\parindent}{0pt}
    \setlength{\parskip}{6pt plus 2pt minus 1pt}}
}{% if KOMA class
  \KOMAoptions{parskip=half}}
\makeatother
\usepackage{color}
\usepackage{fancyvrb}
\newcommand{\VerbBar}{|}
\newcommand{\VERB}{\Verb[commandchars=\\\{\}]}
\DefineVerbatimEnvironment{Highlighting}{Verbatim}{commandchars=\\\{\}}
% Add ',fontsize=\small' for more characters per line
\usepackage{framed}
\definecolor{shadecolor}{RGB}{248,248,248}
\newenvironment{Shaded}{\begin{snugshade}}{\end{snugshade}}
\newcommand{\AlertTok}[1]{\textcolor[rgb]{0.94,0.16,0.16}{#1}}
\newcommand{\AnnotationTok}[1]{\textcolor[rgb]{0.56,0.35,0.01}{\textbf{\textit{#1}}}}
\newcommand{\AttributeTok}[1]{\textcolor[rgb]{0.13,0.29,0.53}{#1}}
\newcommand{\BaseNTok}[1]{\textcolor[rgb]{0.00,0.00,0.81}{#1}}
\newcommand{\BuiltInTok}[1]{#1}
\newcommand{\CharTok}[1]{\textcolor[rgb]{0.31,0.60,0.02}{#1}}
\newcommand{\CommentTok}[1]{\textcolor[rgb]{0.56,0.35,0.01}{\textit{#1}}}
\newcommand{\CommentVarTok}[1]{\textcolor[rgb]{0.56,0.35,0.01}{\textbf{\textit{#1}}}}
\newcommand{\ConstantTok}[1]{\textcolor[rgb]{0.56,0.35,0.01}{#1}}
\newcommand{\ControlFlowTok}[1]{\textcolor[rgb]{0.13,0.29,0.53}{\textbf{#1}}}
\newcommand{\DataTypeTok}[1]{\textcolor[rgb]{0.13,0.29,0.53}{#1}}
\newcommand{\DecValTok}[1]{\textcolor[rgb]{0.00,0.00,0.81}{#1}}
\newcommand{\DocumentationTok}[1]{\textcolor[rgb]{0.56,0.35,0.01}{\textbf{\textit{#1}}}}
\newcommand{\ErrorTok}[1]{\textcolor[rgb]{0.64,0.00,0.00}{\textbf{#1}}}
\newcommand{\ExtensionTok}[1]{#1}
\newcommand{\FloatTok}[1]{\textcolor[rgb]{0.00,0.00,0.81}{#1}}
\newcommand{\FunctionTok}[1]{\textcolor[rgb]{0.13,0.29,0.53}{\textbf{#1}}}
\newcommand{\ImportTok}[1]{#1}
\newcommand{\InformationTok}[1]{\textcolor[rgb]{0.56,0.35,0.01}{\textbf{\textit{#1}}}}
\newcommand{\KeywordTok}[1]{\textcolor[rgb]{0.13,0.29,0.53}{\textbf{#1}}}
\newcommand{\NormalTok}[1]{#1}
\newcommand{\OperatorTok}[1]{\textcolor[rgb]{0.81,0.36,0.00}{\textbf{#1}}}
\newcommand{\OtherTok}[1]{\textcolor[rgb]{0.56,0.35,0.01}{#1}}
\newcommand{\PreprocessorTok}[1]{\textcolor[rgb]{0.56,0.35,0.01}{\textit{#1}}}
\newcommand{\RegionMarkerTok}[1]{#1}
\newcommand{\SpecialCharTok}[1]{\textcolor[rgb]{0.81,0.36,0.00}{\textbf{#1}}}
\newcommand{\SpecialStringTok}[1]{\textcolor[rgb]{0.31,0.60,0.02}{#1}}
\newcommand{\StringTok}[1]{\textcolor[rgb]{0.31,0.60,0.02}{#1}}
\newcommand{\VariableTok}[1]{\textcolor[rgb]{0.00,0.00,0.00}{#1}}
\newcommand{\VerbatimStringTok}[1]{\textcolor[rgb]{0.31,0.60,0.02}{#1}}
\newcommand{\WarningTok}[1]{\textcolor[rgb]{0.56,0.35,0.01}{\textbf{\textit{#1}}}}
\usepackage{graphicx}
\makeatletter
\newsavebox\pandoc@box
\newcommand*\pandocbounded[1]{% scales image to fit in text height/width
  \sbox\pandoc@box{#1}%
  \Gscale@div\@tempa{\textheight}{\dimexpr\ht\pandoc@box+\dp\pandoc@box\relax}%
  \Gscale@div\@tempb{\linewidth}{\wd\pandoc@box}%
  \ifdim\@tempb\p@<\@tempa\p@\let\@tempa\@tempb\fi% select the smaller of both
  \ifdim\@tempa\p@<\p@\scalebox{\@tempa}{\usebox\pandoc@box}%
  \else\usebox{\pandoc@box}%
  \fi%
}
% Set default figure placement to htbp
\def\fps@figure{htbp}
\makeatother
\setlength{\emergencystretch}{3em} % prevent overfull lines
\providecommand{\tightlist}{%
  \setlength{\itemsep}{0pt}\setlength{\parskip}{0pt}}
\usepackage{bookmark}
\IfFileExists{xurl.sty}{\usepackage{xurl}}{} % add URL line breaks if available
\urlstyle{same}
\hypersetup{
  pdftitle={Planejamento fatorial completo com duplicata na área térmica},
  pdfauthor={Andres Stiven Carreño Jerez},
  hidelinks,
  pdfcreator={LaTeX via pandoc}}

\title{Planejamento fatorial completo com duplicata na área térmica}
\author{Andres Stiven Carreño Jerez}
\date{2026-05-26}

\begin{document}
\maketitle

\section{Construção Fatorial completo 2² con duplicata
experimental}\label{construuxe7uxe3o-fatorial-completo-2uxb2-con-duplicata-experimental}

\begin{Shaded}
\begin{Highlighting}[]
\FunctionTok{library}\NormalTok{(FrF2)}
\end{Highlighting}
\end{Shaded}

\begin{verbatim}
## Cargando paquete requerido: DoE.base
\end{verbatim}

\begin{verbatim}
## Cargando paquete requerido: grid
\end{verbatim}

\begin{verbatim}
## Cargando paquete requerido: conf.design
\end{verbatim}

\begin{verbatim}
## Registered S3 method overwritten by 'DoE.base':
##   method           from       
##   factorize.factor conf.design
\end{verbatim}

\begin{verbatim}
## 
## Adjuntando el paquete: 'DoE.base'
\end{verbatim}

\begin{verbatim}
## The following objects are masked from 'package:stats':
## 
##     aov, lm
\end{verbatim}

\begin{verbatim}
## The following object is masked from 'package:graphics':
## 
##     plot.design
\end{verbatim}

\begin{verbatim}
## The following object is masked from 'package:base':
## 
##     lengths
\end{verbatim}

\begin{Shaded}
\begin{Highlighting}[]
\NormalTok{dados1}\OtherTok{=}\FunctionTok{data.frame}\NormalTok{(}\AttributeTok{A =} \FunctionTok{c}\NormalTok{(}\SpecialCharTok{{-}}\DecValTok{1}\NormalTok{,}\SpecialCharTok{{-}}\DecValTok{1}\NormalTok{,}\DecValTok{1}\NormalTok{,}\DecValTok{1}\NormalTok{,}\SpecialCharTok{{-}}\DecValTok{1}\NormalTok{,}\SpecialCharTok{{-}}\DecValTok{1}\NormalTok{,}\DecValTok{1}\NormalTok{,}\DecValTok{1}\NormalTok{),}
                    \AttributeTok{B =} \FunctionTok{c}\NormalTok{(}\SpecialCharTok{{-}}\DecValTok{1}\NormalTok{,}\SpecialCharTok{{-}}\DecValTok{1}\NormalTok{,}\SpecialCharTok{{-}}\DecValTok{1}\NormalTok{,}\SpecialCharTok{{-}}\DecValTok{1}\NormalTok{,}\DecValTok{1}\NormalTok{,}\DecValTok{1}\NormalTok{,}\DecValTok{1}\NormalTok{,}\DecValTok{1}\NormalTok{), }
                  \AttributeTok{rep =} \FunctionTok{c}\NormalTok{(}\DecValTok{1}\NormalTok{,}\DecValTok{2}\NormalTok{,}\DecValTok{1}\NormalTok{,}\DecValTok{2}\NormalTok{,}\DecValTok{1}\NormalTok{,}\DecValTok{2}\NormalTok{,}\DecValTok{1}\NormalTok{,}\DecValTok{2}\NormalTok{),}
                    \AttributeTok{k =} \FunctionTok{c}\NormalTok{(}\FloatTok{0.31}\NormalTok{,}\FloatTok{0.33}\NormalTok{,}\FloatTok{0.58}\NormalTok{,}\FloatTok{0.61}\NormalTok{,}\FloatTok{0.35}\NormalTok{,}\FloatTok{0.36}\NormalTok{,}\FloatTok{0.72}\NormalTok{,}\FloatTok{0.75}\NormalTok{))}

\NormalTok{dados1}
\end{Highlighting}
\end{Shaded}

\begin{verbatim}
##    A  B rep    k
## 1 -1 -1   1 0.31
## 2 -1 -1   2 0.33
## 3  1 -1   1 0.58
## 4  1 -1   2 0.61
## 5 -1  1   1 0.35
## 6 -1  1   2 0.36
## 7  1  1   1 0.72
## 8  1  1   2 0.75
\end{verbatim}

\section{Médias e desvios por
condição}\label{muxe9dias-e-desvios-por-condiuxe7uxe3o}

\begin{Shaded}
\begin{Highlighting}[]
\NormalTok{condicoes}\OtherTok{=}\FunctionTok{aggregate}\NormalTok{(k }\SpecialCharTok{\textasciitilde{}}\NormalTok{ A }\SpecialCharTok{+}\NormalTok{ B, }\AttributeTok{data =}\NormalTok{ dados1,}
                       \AttributeTok{FUN =} \ControlFlowTok{function}\NormalTok{(x) }\FunctionTok{c}\NormalTok{(}\AttributeTok{media =} \FunctionTok{mean}\NormalTok{(x), }\AttributeTok{dp =} \FunctionTok{sd}\NormalTok{(x)))}
\FunctionTok{cat}\NormalTok{(}\StringTok{"}\SpecialCharTok{\textbackslash{}n}\StringTok{=== MÉDIA E DP POR CONDIÇÃO ===}\SpecialCharTok{\textbackslash{}n}\StringTok{"}\NormalTok{)}
\end{Highlighting}
\end{Shaded}

\begin{verbatim}
## 
## === MÉDIA E DP POR CONDIÇÃO ===
\end{verbatim}

\begin{Shaded}
\begin{Highlighting}[]
\FunctionTok{print}\NormalTok{(condicoes)}
\end{Highlighting}
\end{Shaded}

\begin{verbatim}
##    A  B     k.media        k.dp
## 1 -1 -1 0.320000000 0.014142136
## 2  1 -1 0.595000000 0.021213203
## 3 -1  1 0.355000000 0.007071068
## 4  1  1 0.735000000 0.021213203
\end{verbatim}

\section{Calcular efeitos}\label{calcular-efeitos}

\begin{Shaded}
\begin{Highlighting}[]
\NormalTok{y}\OtherTok{=}\NormalTok{dados1}\SpecialCharTok{$}\NormalTok{k}
\NormalTok{sA}\OtherTok{=}\NormalTok{dados1}\SpecialCharTok{$}\NormalTok{A}
\NormalTok{sB}\OtherTok{=}\NormalTok{dados1}\SpecialCharTok{$}\NormalTok{B}
\NormalTok{sAB}\OtherTok{=}\NormalTok{sA }\SpecialCharTok{*}\NormalTok{ sB}

\NormalTok{efecto }\OtherTok{\textless{}{-}} \ControlFlowTok{function}\NormalTok{(s, y) }\FunctionTok{sum}\NormalTok{(s }\SpecialCharTok{*}\NormalTok{ y) }\SpecialCharTok{/}\NormalTok{ (}\FunctionTok{length}\NormalTok{(y)}\SpecialCharTok{/}\DecValTok{2}\NormalTok{)}

\FunctionTok{cat}\NormalTok{(}\StringTok{"}\SpecialCharTok{\textbackslash{}n}\StringTok{=== EFEITOS ===}\SpecialCharTok{\textbackslash{}n}\StringTok{"}\NormalTok{)}
\end{Highlighting}
\end{Shaded}

\begin{verbatim}
## 
## === EFEITOS ===
\end{verbatim}

\begin{Shaded}
\begin{Highlighting}[]
\FunctionTok{cat}\NormalTok{(}\FunctionTok{sprintf}\NormalTok{(}\StringTok{"Efeito A  (Grafite)   : \%+.5f W/(m·K)}\SpecialCharTok{\textbackslash{}n}\StringTok{"}\NormalTok{, }\FunctionTok{efecto}\NormalTok{(sA,  y)))}
\end{Highlighting}
\end{Shaded}

\begin{verbatim}
## Efeito A  (Grafite)   : +0.32750 W/(m·K)
\end{verbatim}

\begin{Shaded}
\begin{Highlighting}[]
\FunctionTok{cat}\NormalTok{(}\FunctionTok{sprintf}\NormalTok{(}\StringTok{"Efeito B  (Temperatura): \%+.5f W/(m·K)}\SpecialCharTok{\textbackslash{}n}\StringTok{"}\NormalTok{, }\FunctionTok{efecto}\NormalTok{(sB,  y)))}
\end{Highlighting}
\end{Shaded}

\begin{verbatim}
## Efeito B  (Temperatura): +0.08750 W/(m·K)
\end{verbatim}

\begin{Shaded}
\begin{Highlighting}[]
\FunctionTok{cat}\NormalTok{(}\FunctionTok{sprintf}\NormalTok{(}\StringTok{"Efeito AB              : \%+.5f W/(m·K)}\SpecialCharTok{\textbackslash{}n}\StringTok{"}\NormalTok{, }\FunctionTok{efecto}\NormalTok{(sAB, y)))}
\end{Highlighting}
\end{Shaded}

\begin{verbatim}
## Efeito AB              : +0.05250 W/(m·K)
\end{verbatim}

\section{ANOVA}\label{anova}

\begin{Shaded}
\begin{Highlighting}[]
\NormalTok{dados1}\SpecialCharTok{$}\NormalTok{A }\OtherTok{\textless{}{-}} \FunctionTok{as.factor}\NormalTok{(dados1}\SpecialCharTok{$}\NormalTok{A)}
\NormalTok{dados1}\SpecialCharTok{$}\NormalTok{B }\OtherTok{\textless{}{-}} \FunctionTok{as.factor}\NormalTok{(dados1}\SpecialCharTok{$}\NormalTok{B)}

\NormalTok{modelo7 }\OtherTok{\textless{}{-}} \FunctionTok{lm}\NormalTok{(k }\SpecialCharTok{\textasciitilde{}}\NormalTok{ A }\SpecialCharTok{*}\NormalTok{ B, }\AttributeTok{data =}\NormalTok{ dados1)}

\FunctionTok{print}\NormalTok{(}\FunctionTok{anova}\NormalTok{(modelo7))}
\end{Highlighting}
\end{Shaded}

\begin{verbatim}
## Analysis of Variance Table
## 
## Response: k
##           Df   Sum Sq  Mean Sq F value    Pr(>F)    
## A          1 0.214513 0.214513 746.130 1.068e-05 ***
## B          1 0.015313 0.015313  53.261  0.001874 ** 
## A:B        1 0.005512 0.005512  19.174  0.011886 *  
## Residuals  4 0.001150 0.000288                      
## ---
## Signif. codes:  0 '***' 0.001 '**' 0.01 '*' 0.05 '.' 0.1 ' ' 1
\end{verbatim}

\begin{Shaded}
\begin{Highlighting}[]
\FunctionTok{print}\NormalTok{(}\FunctionTok{summary}\NormalTok{(modelo7))}
\end{Highlighting}
\end{Shaded}

\begin{verbatim}
## 
## Call:
## lm.default(formula = k ~ A * B, data = dados1)
## 
## Residuals:
##      1      2      3      4      5      6      7      8 
## -0.010  0.010 -0.015  0.015 -0.005  0.005 -0.015  0.015 
## 
## Coefficients:
##             Estimate Std. Error t value Pr(>|t|)    
## (Intercept)  0.32000    0.01199  26.690 1.17e-05 ***
## A1           0.27500    0.01696  16.219 8.46e-05 ***
## B1           0.03500    0.01696   2.064   0.1079    
## A1:B1        0.10500    0.02398   4.379   0.0119 *  
## ---
## Signif. codes:  0 '***' 0.001 '**' 0.01 '*' 0.05 '.' 0.1 ' ' 1
## 
## Residual standard error: 0.01696 on 4 degrees of freedom
## Multiple R-squared:  0.9951, Adjusted R-squared:  0.9915 
## F-statistic: 272.9 on 3 and 4 DF,  p-value: 4.427e-05
\end{verbatim}

\section{Grafo Efe}\label{grafo-efe}

\begin{Shaded}
\begin{Highlighting}[]
\NormalTok{plano7}\OtherTok{=}\FunctionTok{FrF2}\NormalTok{(}\AttributeTok{nruns    =} \DecValTok{4}\NormalTok{,}\AttributeTok{nfactors =} \DecValTok{2}\NormalTok{,}\AttributeTok{factor.names =} \FunctionTok{list}\NormalTok{(}\AttributeTok{A =} \FunctionTok{c}\NormalTok{(}\SpecialCharTok{{-}}\DecValTok{1}\NormalTok{, }\DecValTok{1}\NormalTok{),}\AttributeTok{B =} \FunctionTok{c}\NormalTok{(}\SpecialCharTok{{-}}\DecValTok{1}\NormalTok{, }\DecValTok{1}\NormalTok{)),}\AttributeTok{randomize =} \ConstantTok{FALSE}
\NormalTok{)}
\end{Highlighting}
\end{Shaded}

\begin{verbatim}
## creating full factorial with 4 runs ...
\end{verbatim}

\begin{Shaded}
\begin{Highlighting}[]
\CommentTok{\# ✅ Agregar respuesta con add.response() — NO con $}
\NormalTok{k\_medias}\OtherTok{=}\FunctionTok{c}\NormalTok{(}\FloatTok{0.320}\NormalTok{, }\FloatTok{0.595}\NormalTok{, }\FloatTok{0.355}\NormalTok{, }\FloatTok{0.735}\NormalTok{)}
\NormalTok{plano7}\OtherTok{=}\FunctionTok{add.response}\NormalTok{(plano7, k\_medias)}

\CommentTok{\# Verificar que quedó bien}
\FunctionTok{print}\NormalTok{(plano7)}
\end{Highlighting}
\end{Shaded}

\begin{verbatim}
##    A  B k_medias
## 1 -1 -1    0.320
## 2  1 -1    0.595
## 3 -1  1    0.355
## 4  1  1    0.735
## class=design, type= full factorial
\end{verbatim}

\begin{Shaded}
\begin{Highlighting}[]
\CommentTok{\# Gráfico de efectos principales}
\FunctionTok{MEPlot}\NormalTok{(plano7,}
       \AttributeTok{response =} \StringTok{"k\_medias"}\NormalTok{,}
       \AttributeTok{main     =} \StringTok{"Efeitos Principais – Condutividade Térmica k"}\NormalTok{)}
\end{Highlighting}
\end{Shaded}

\pandocbounded{\includegraphics[keepaspectratio]{7._files/7._files/figure-latex/unnamed-chunk-5-1.pdf}}

\begin{Shaded}
\begin{Highlighting}[]
\CommentTok{\# Gráfico de interacción}
\FunctionTok{IAPlot}\NormalTok{(plano7,}
       \AttributeTok{response =} \StringTok{"k\_medias"}\NormalTok{,}
       \AttributeTok{main     =} \StringTok{"Interação A×B – Grafite vs. Temperatura"}\NormalTok{)}
\end{Highlighting}
\end{Shaded}

\pandocbounded{\includegraphics[keepaspectratio]{7._files/7._files/figure-latex/unnamed-chunk-5-2.pdf}}

\section{Grafo Efe}\label{grafo-efe-1}

El factor más importante fue: ★ A -- Teor de grafite expandido (efecto =
+0,3275 W·m⁻¹·K⁻¹)

Su efecto POSITIVO indica que pasar del nivel bajo al nivel alto AUMENTA
fuertemente la conductividad: - Nivel bajo = 5\% de grafite expandido -
Nivel alto = 15\% de grafite expandido → Mayor concentración de grafite
maximiza la conducción térmica del compósito.

═══════════════════════════════════════════════════════════ EFECTOS
PRINCIPALES ═══════════════════════════════════════════════════════════
A (Teor grafite) : +0,3275 W·m⁻¹·K⁻¹ → DOMINANTE ★ B (Temp. prensagem) :
+0,0875 W·m⁻¹·K⁻¹ → moderado A×B (interacción) : +0,0525 W·m⁻¹·K⁻¹ →
significativa

═══════════════════════════════════════════════════════════ ANOVA --
SIGNIFICANCIA ESTADÍSTICA
═══════════════════════════════════════════════════════════ Fuente F
calculado F crítico(5\%) Conclusión
────────────────────────────────────────────────── A 746,1 7,71
SIGNIFICATIVO ✅ B 53,3 7,71 SIGNIFICATIVO ✅ A×B 19,2 7,71
SIGNIFICATIVO ✅ ────────────────────────────────────────────────── Los
TRES efectos son estadísticamente significativos.

═══════════════════════════════════════════════════════════ INTERACCIÓN
A×B --- SINERGIA ENTRE FACTORES
═══════════════════════════════════════════════════════════ Con bajo
grafite (5\%): B sube k apenas +0,035 W·m⁻¹·K⁻¹

Con alto grafite (15\%): B sube k en +0,140 W·m⁻¹·K⁻¹

→ La temperatura de prensagem ayuda MUCHO MÁS cuando hay alta
concentración de grafite. → Los dos factores actúan en SINERGIA ---
juntos producen un efecto mayor que la suma individual.

═══════════════════════════════════════════════════════════
REPETIBILIDAD DEL EXPERIMENTO
═══════════════════════════════════════════════════════════ Condición k̄
(W·m⁻¹·K⁻¹) DP ───────────────────────────────────────── A=5\%, T=170°C
0,320 0,0141 A=15\%, T=170°C 0,595 0,0212 A=5\%, T=190°C 0,355 0,0071
A=15\%, T=190°C 0,735 0,0212 ─────────────────────────────────────────
Todos los DP son pequeños → alta repetibilidad ✅ El experimento es
confiable y reproducible.

═══════════════════════════════════════════════════════════ MEJOR
CONDICIÓN EXPERIMENTAL
═══════════════════════════════════════════════════════════ Para
MAXIMIZAR la conductividad térmica:

→ ENSAYO condición 4: - Teor de grafite : 15\% (nivel alto) - Temp. de
prensagem : 190°C (nivel alto) - k média : 0,735 W·m⁻¹·K⁻¹ ← mayor valor
- DP : 0,021 ← repetible ✅ - CV : 2,9\% ← excelente ✅

Esta condición aprovecha la SINERGIA A×B: el grafite forma redes
conductoras (percolación) y la alta temperatura mejora el contacto entre
partículas durante la prensagem.

\end{document}
